% ============================================================
% A Learned Proxy for the Shadow of Intervention
% ============================================================

\section{A Learned Proxy for the Shadow of Intervention}
\label{sec:constructing}

Operationalizing the shadow requires a measure of the expected intervention
environment at the moment a potential rebel group weighs the costs and benefits
of open conflict.  The construction proceeds in two stages: a directed-dyad
classifier estimates intervention probabilities for every potential
intervener, and a country-year aggregation sums them under a Nash
fixed-point condition ensuring self-consistency.

\subsection{Outcome Variable}

The Stage~1 training labels draw primarily from \citet{regan2000},
whose dataset records third-party military, economic, and diplomatic
interventions in civil conflicts from 1944 to 1999.  I extend Regan's coding
through 2014 by hand-coding unambiguous military interventions in post-1999 COW
onset country-years.  The coding threshold requires deployment of military
personnel---troops, combat advisors, or directed proxy forces attributable to a
specific state---on behalf of one side; arms transfers, economic aid, and
sanctuary alone are insufficient.  UN peacekeeping missions are excluded.  The
post-1999 extension adds 31 directed dyad-year records across 14 host countries,
yielding a combined intervention table of 455 records (1944--2014); the full
coding with source notes appears in Online Appendix~\ref{app:coding}.

The civil war universe is drawn from the Correlates of War Intra-State War
Dataset (v5.1) \citep{dixon2016}, restricting to wars fought for central control
(type~4) or local/secessionist issues (type~5).  The COW threshold
($\geq$1,000 battle deaths) is theoretically appropriate: the mechanism requires
rebels organized enough to form rational expectations and conflicts large
enough that outside military involvement is a realistic prospect.\footnote{%
  At lower thresholds---such as the UCDP/PRIO $\geq$25-death criterion---third-party
  military intervention becomes vanishingly rare, swamping any predictive signal.}
This yields 153 distinct onset country-years between 1946 and 1999 and 191
through 2014.  I restrict attention to \textit{military} interventions, coded as
\textit{government-biased}~(1) or \textit{opposition-biased}~(2); corrections and
deduplication are documented in Online Appendix~\ref{app:coding}.

For each directed dyad $(B \rightarrow A)$ active in a civil war onset year, the
training label is: 0 (no military intervention), 1 (government-biased), or 2
(opposition-biased).  Non-onset dyad-years are excluded from the training set.
This produces 254 intervention-coded onset observations---150 government-biased
and 104 opposition-biased---across 51 host countries and 58 unique intervening
states.  The five permanent Security Council members
account for 110 of these events (43~percent); the remaining
states---Cold War proxies, neighbors, regional patrons, and coethnics---account
for the majority (Table~\ref{tab:interveners}).

Training exclusively on onset country-years is not a convenience---it is
the only option.  Intervention in a civil war is undefined absent a civil war;
in non-onset years the outcome is not ``zero intervention'' but rather
unobserved.  The consequence is that the training sample is selected on
precisely the outcome that the Stage~2 analysis tries to explain: if the
shadow deters onset, then high-shadow country-years are systematically
underrepresented among onsets, and the classifier learns the
feature--intervention mapping from an endogenously filtered sample.  There
is no design that avoids this.  The defense rests on what the classifier
is asked to learn.  The training target is the conditional relationship
between \textit{dyad-level} features---alliance portfolios, capability ratios,
colonial ties, geographic proximity, regime similarity---and the direction
of military intervention, given that a war has started.  These bilateral
characteristics are plausible determinants of how state $B$ behaves toward
state $A$ in a conflict regardless of what brought that conflict about.
The fixed-point iteration then
handles the extrapolation: it asks what the equilibrium intervention
environment \textit{would} look like in any country-year, including those where
onset was never observed, without requiring that the training sample be
representative of the full population.

\begin{table}[t]
  \centering
  \begin{tabular}{@{}lrrr@{}} 
    \toprule
    State & Gov. & Opp. & Total \\
    \midrule
    United States\textsuperscript{\dag}   & 31 & 15 & 46 \\
    USSR / Russia\textsuperscript{\dag}   & 16 &  7 & 23 \\
    United Kingdom\textsuperscript{\dag}  &  9 &  8 & 17 \\
    France\textsuperscript{\dag}          & 11 &  6 & 17 \\
    Cuba                                  &  7 &  1 &  8 \\
    China\textsuperscript{\dag}           &  2 &  5 &  7 \\
    Libya                                 &  1 &  6 &  7 \\
    Iran                                  &  2 &  5 &  7 \\
    Somalia                               &  0 &  5 &  5 \\
    Syria                                 &  3 &  2 &  5 \\
    India                                 &  4 &  1 &  5 \\
    Uganda                                &  3 &  2 &  5 \\
    \bottomrule
  \end{tabular}
  \caption{Most frequent military interveners (1946--2014); states with
    fewer than five interventions omitted.
    \textsuperscript{\dag}~Permanent member of the UN Security Council (P5).
    Government-biased (Gov.)\ and opposition-biased (Opp.)\ interventions
    are coded from the Regan \texttt{target} variable and post-1999 extension.}
  \label{tab:interveners}
\end{table}

\subsection{Feature Set}

The classifier draws on 105 directed-dyad-year features organized
around four substantive dimensions (full list in Online Appendix~\ref{app:varlist}).
Many variables enter twice---once for the potential intervener and
once for the host---so that the classifier can learn asymmetric effects.
Missing data are multiply imputed using miceforest (five country-year
$\times$ five undirected-dyad imputations, yielding 25 complete datasets;
Online Appendix~\ref{app:imputation}).

\begin{itemize}
\item \textbf{State characteristics of both the potential intervener and the host.}
  Material capabilities (CINC scores and components, GDP, population) from COW
  v6 \citep{singer1987} and V-Dem v15 \citep{vdem2025}; regime type (Polity~II and lags,
  liberal democracy index, instability); power status (COW major power, P5
  membership); and host-side conflict history, ethnic fractionalization
  \citep{fearon2003}, ethnic exclusion \citep{cederman2010}, oil wealth, and terrain.

\item \textbf{Bilateral ties.}  Alliance commitments \citep{gibler2009}, bilateral trade,
  geographic proximity (log capital distance, contiguity \citealp{stinnett2002}),
  colonial history, territorial disputes \citep{huth2003}, expected dispute outcomes
  \citep{carroll2019}, and peace quality \citep{diehl2019}.

\item \textbf{Foreign policy alignment.}  UN General Assembly ideal points
  \citep{bailey2017} for each state and the absolute distance between them;
  shared IGO membership from COW Dyadic IGO v3 \citep{pevehouse2020}.

\item \textbf{Spatial context.}  Ten spatial lag variables constructed from a
  row-normalized polity-similarity weight matrix capture the weighted behavior
  of other potential interveners in the same host-year: the fraction
  coded government- or opposition-biased, US and Soviet intervention
  posture, and four superpower interaction terms \citep{gent2007,salehyan2007}
  (Online Appendix~\ref{app:spatial}).
\end{itemize}

Before training, the 115 features (95 dyadic covariates + 20 spatial lags) are
standardized and reduced via principal components analysis, retaining components
to 90\% of cumulative variance (59--61~components across all 25 imputation
datasets).

\subsection{Classifier Design}

The multinomial outcome (0 = no intervention, 1 = government-biased,
2 = opposition-biased) is modeled using a super-learner ensemble
\citep{vanderlaan2007}.

Nine component classifiers spanning three model families are trained
within a ten-fold cross-validation scheme: two tree ensembles (a random
forest \citealp{breiman2001} and two histogram-based gradient boosting
specifications at different learning rates), four logistic regression variants (ridge,
elastic-net, lasso, and unpenalized multinomial), and two multi-layer
perceptrons (shallow and deeper architectures).  The unpenalized
multinomial logistic regression---a softmax best-response function---directly
nests the functional-form assumption of structural game-theoretic models of
intervention \citep{gibilisco2022}; its stacking weight is determined by
out-of-fold predictive performance, so whether that assumption fits the data
is answered empirically rather than imposed by construction.  Each
classifier produces a vector of class probabilities for every
training observation.  The super learner combines these using
non-negative least squares (NNLS) weights chosen to minimize held-out
multinomial log-loss, so that the ensemble adaptively up-weights
component classifiers that predict well on out-of-fold observations.

All reported performance statistics are based on out-of-fold
predictions to guard against overfitting---following \citeauthor{wang2018}'s
(\citeyear{wang2018}) demonstration that
evaluation of class-imbalanced political-event models on training data
produces severely inflated performance estimates \citep{muchlinski2016}.
The primary performance metric is the proportional reduction in log-loss
relative to the class-frequency null model (PRL); the area under the ROC
curve (AUC) is also reported.

\paragraph{Spatial lags and Nash fixed-point iteration.}
The broader motivation for
the two-stage design is that intervention decisions are endogenous to the
conflict environment: interveners select into conflicts strategically, and
a single-equation model that treats realized intervention as a covariate in
an onset regression conflates the causal effect with the selection process
that generated it \citep{signorino2002,gent2008}.  Qualitative evidence confirms
the problem directly: rebel groups in Bosnia and Kosovo explicitly
conditioned their onset decisions on anticipated outside support, so
realized intervention is partly a consequence of expected intervention
rather than an independent cause \citep{kuperman2008}.  Using predicted
intervention expectations rather than realized intervention is the
appropriate correction.

The spatial context variables described in Appendix~\ref{app:spatial} depend on the
intervention choices of \textit{other} potential interveners---quantities that are
themselves to be predicted.  A documented pattern of countervailing
intervention makes these dependencies quantitatively important: when one
state intervenes on the government side, ideologically rival states face a
selective incentive to support the opposition, and vice versa \citep{salehyan2011,findleyteo2006}.
\citet{aydinregan2011} confirm that these
network effects are consequential: opposing interventions nearly double civil
war duration, while same-side interventions shorten wars only when the
intervening states share similar foreign policy preferences---evidence that
the strategic alignment among interveners, not merely their count, carries
the signal.  This creates a concrete measurement problem.
Without an explicit consistency requirement, spatial lags are computed
from \textit{observed} training-data intervention choices, which are not the
model's own outputs.  Applied out-of-sample---to country-years unlike
the training-period equilibrium---the predictions embed a deterministic
discrepancy between lag inputs and probability outputs whose sign and
magnitude depend on how far the out-of-sample world departs from
training-period behavior.  This is systematic error, not sampling noise,
and in a measure designed to capture counterfactual intervention
environments the bias could be large.

The fix is to require that the predictions be self-consistent inputs to
themselves:

\begin{equation}
\sigma^*(B, A, t) = \hat{M}(X(B, A, t, \sigma^*)) \quad \forall\, (B, A, t),
\label{eq:fp}
\end{equation}

\noindent where $\hat{M}$ is the trained super-learner classifier.  At $\sigma^*$,
the spatial lags are derived from the model's own equilibrium predictions,
not from historical observations.
The deterministic discrepancy disappears by construction.

This self-consistency condition is also the Nash equilibrium condition
for a natural class of games.  The primitive object is the best-response
correspondence, not the game itself: a game specifies a best-response
correspondence \textit{derivatively}, via utility maximization; $\hat{M}$ specifies
one \textit{directly}, from the data.  Nash equilibrium is a fixed point of
the best-response correspondence and does not require identifying the
payoff functions that generate it.  Games can be partitioned into
equivalence classes under the relation ``generates the same best-response
correspondence''; $\sigma^*$ is the Nash equilibrium of every game in the
class whose correspondence $\hat{M}$ estimates.\footnote{%
  The set of rationalizing payoff functions for any interior $\sigma^*$
  is non-empty---choose utilities that make each player indifferent among
  actions assigned positive probability.  What is identified is the
  correspondence; the payoffs that generate it are not.  A natural
  parametric fiber of these classes is the linear-utility family
  $u_i(a, \sigma_{-i}) = \alpha \cdot f_i(a) + \beta \cdot W_i(a, \sigma_{-i})$,
  where $f_i$ captures state-specific intervention incentives and $W_i$
  the spatial interaction term.  The best-response from this family is a
  softmax---the functional form of the multinomial logit component of
  the super-learner---so the parametric case is explicitly represented
  in the ensemble while the broader non-parametric mixture accommodates
  departures from it.  The iterative procedure that solves Equation~\eqref{eq:fp}
  is the nonparametric analogue of the nested pseudo-likelihood (NPL) algorithm
  of \citet{aguirregabiriamira2007}: the fitted
  super-learner replaces the parametric best-response function,
  and the fixed-point iteration updates beliefs (spatial lags) until
  self-consistency, as in A\&M's two-step PML.  In the parametric case,
  \citet{crismancox2021} show that this class of
  iterative estimators substantially outperforms direct maximum likelihood
  estimation of strategic games under equilibrium multiplicity.
  Uniqueness of the fixed point is not guaranteed in general, but
  convergence is rapid (two to three passes across all 25 imputation
  draws), consistent with a contractive mapping.}
If the training data reflect equilibrium play and the spatial lags are
sufficient statistics for the strategically relevant information in
$\sigma_{-i}$, then $\hat{M}$ is a non-parametric estimate of the true
best-response correspondence $\mathrm{BR}^*$ and $\sigma^*$ estimates its fixed
point, without specifying which game in the equivalence class is the
true data-generating process.  The statistical
and game-theoretic justifications therefore coincide: Equation~\eqref{eq:fp} corrects a
concrete measurement error and simultaneously identifies the equilibrium
of the underlying intervention game, whichever member of the equivalence
class that game happens to be.

Equation~\eqref{eq:fp} has no closed form but is solved by fixed-point iteration.
The iteration proceeds in two stages.  In the \textit{training stage}, the
classifier is first fit using observed spatial lags; out-of-fold predictions
then replace the observed intervention indicators to recompute spatial lags;
and the classifier is refit on the updated lags.  This loop repeats until
convergence, typically within two to three passes.  In the \textit{prediction
stage}, the trained classifier is held fixed and only the spatial lags are
updated: the model predicts for every directed dyad-year (onset and non-onset
alike), spatial lags are recomputed from those predictions, and the model
re-predicts until convergence.  Applying the fixed-point universally---not just to onset rows---ensures
that the shadow reflects the counterfactual equilibrium in any country-year,
including those where no conflict was observed.

\subsection{Predictive Performance}

Table~\ref{tab:stage1-perf} reports out-of-fold performance for the super-learner
and two selected components: the lasso logistic regression (the best single
component by standalone PRL) and the random forest (which despite lower standalone
PRL receives the second-highest NNLS weight at 38.6\%, reflecting complementary
predictive information).  Full results for all nine component learners appear in
Appendix~\ref{app:sl}, Table~\ref{tab:sl-components}.  All statistics average
across the 25 complete datasets, with standard deviations across imputation draws
in parentheses.

\begin{table}[t]
  \centering
  \begin{tabular}{@{}lrrr@{}}
    \toprule
    Model & PRL & AUC (macro) & Log-loss \\
    \midrule
    Super-learner       & 47.3\% (1.0\%) & 0.969 (0.003) & 0.033 (0.001) \\
    MLP (100, 50)       & 33.4\% (4.4\%) & 0.940 (0.014) & 0.041 (0.003) \\
    Lasso logit         & 35.1\% (1.0\%) & 0.957 (0.003) & 0.040 (0.001) \\
    \bottomrule
  \end{tabular}
  \caption{Selected Stage~1 classifier performance (out-of-fold, 25 imputation
    draws).  The MLP is the highest-weighted component (44.5\% of the ensemble);
    lasso logit is the best parametric model.  Standard deviations across draws
    in parentheses.  The class-frequency null has PRL~$=$~0\% by construction,
    AUC~$=$~0.500, and log-loss~$=$~0.062.  Full results for all nine component
    learners appear in Table~\ref{tab:sl-components}.}
  \label{tab:stage1-perf}
\end{table}

The ensemble achieves a PRL of 47.3\% and a macro-average AUC of 0.969.
A PRL near 47\% indicates that the classifier recovers nearly half the
information in intervention outcomes that the class-frequency null
discards---a substantial figure for a three-class rare-event problem
in which the modal outcome (no intervention) accounts for over 99\% of
observations.  Interventions account for barely one percent of
directed-dyad-year observations, so the predictive task is a
needle-in-a-haystack problem.  A natural question is whether the
classifier earns its performance by finding the needles---correctly
identifying the rare dyad-years where intervention occurs---or by
calibrating the hay, sharpening predictions for the overwhelming
non-intervention majority.  The distinction is consequential: a model
that only improved its calibration of non-events would replicate what
any onset regression without shadow variables already does implicitly,
sending all intervention probabilities to their unconditional near-zero
rate.  What the shadow measure requires is discrimination among rare
events.  Decomposing the log-loss reduction by true class confirms that
the classifier delivers exactly this: 94 percent of the PRL originates
from the 1.0 percent of observations where intervention actually
occurred, with only 6 percent from improved calibration of
non-events.

The result also speaks directly to the rational expectations assumption
embedded in the theoretical model.  The classifier is trained entirely on
publicly observable state characteristics---material capabilities, alliance
portfolios, regime type, colonial history, foreign policy alignment.  If an
estimator operating on this strict subset of information recovers 47\% of
the predictive content in realized intervention outcomes, then intervention
is, to a substantial degree, foreseeable from information available to any
careful observer.  Under rational expectations, agents exploit available
information at least as well as the analyst's model \citep{muth1961}; rebel
leaderships can additionally draw on private intelligence, direct
communications with potential patron states, diaspora networks, and
on-the-ground assessments of great power intentions that no dataset records.
The 47\% figure is accordingly a lower bound on how knowable the
intervention environment is from the rebel's perspective, not a ceiling.

The performance gap between the ensemble and the best single logistic
specification---47.3\% vs 35.2\% PRL on an identical feature set---reveals
nonlinear and interactive structure that parametric models cannot fully recover.
\citet{moruccispirling2024} show that in typical political science applications,
complex models rarely outperform logistic regression because curated feature
sets lack exploitable nonlinear structure; the intervention prediction task
is an exception, precisely because the broad feature space is designed for
measurement rather than theory testing.  The NNLS
stacking procedure quantifies this directly: it assigns near-zero weight to all
four logistic specifications (ridge, lasso, elastic-net, unpenalized
multinomial), concentrating ensemble mass on the deep multi-layer perceptron
(44.5\%) and the random forest (38.6\%), with a gradient-boosted ensemble
taking 11.1\% (Appendix~\ref{app:sl}, Table~\ref{tab:sl-components}).
Notably, the logistic family achieves the highest standalone PRL among
individual learners yet receives negligible ensemble weight, indicating that
the nonlinear learners subsume their predictive content while adding
complementary information the logistic models miss.
Because the logistic family is
explicitly included in the candidate library and explicitly down-weighted by
the data, the ensemble's superiority is not an artifact of architecture choice.

Because Stage~1 probabilities are subsequently summed across approximately
190 potential interveners, systematic miscalibration compounds; reliability
diagrams confirming calibration for each outcome class appear in Appendix~\ref{app:sl}.
These diagnostics follow \citeauthor{knoxlucascho2022}'s (\citeyear{knoxlucascho2022}) recommendation to ``always
assess and report measurement performance'' (their \S5.4): the PRL and AUC
quantify predictive accuracy, and the reliability diagrams verify the
calibration assumptions on which the aggregation to country-year
expectations depends.

\subsection{Aggregating to Country-Year Intervention Expectations}

The Stage~1 classifier produces, for each directed dyad $(B \rightarrow A)$
in a civil war onset country-year, probability estimates
$\hat{\sigma}_B(1)$ and $\hat{\sigma}_B(2)$ for government- and
opposition-biased intervention, respectively.  To reduce noise from
potential interveners to whom the classifier assigns negligible
probability, I apply a cutpoint filter: intervener $B$ contributes to
the aggregate only if $\hat{\sigma}_B(k) \geq \tau$ for $k \in \{1, 2\}$.
The threshold $\tau$ is treated as a tuning parameter and selected by
out-of-fold performance in the Stage~2 onset models; in practice,
$\tau = 0.001$ performs best, excluding only the very least likely
interveners.  The threshold is permissive by design: the PRL
decomposition above shows that the classifier's information content
concentrates overwhelmingly in the rare dyads it flags as plausible
interveners, so the cutpoint discards noise from the long tail of
near-zero predictions while preserving essentially all of the signal.

The aggregate intervention-expectation variables for host country $A$ in
year $t$ are

\begin{equation*}
E_A^G(t) = \sum_{B:\, \hat{\sigma}_B(1) \geq \tau} \hat{\sigma}_B(1), \quad
E_A^O(t) = \sum_{B:\, \hat{\sigma}_B(2) \geq \tau} \hat{\sigma}_B(2).
\end{equation*}

Both sums are mechanically larger in later years as the system expands.
I include a year trend in all second-stage regressions to absorb this
secular growth.  The sums are also right-skewed; following
\citet{burbidge1988}, I apply the inverse hyperbolic sine (asinh)
transformation before entering them into the onset models.

\subsection{Properties of the Shadow Measure}
\label{sec:shadow-properties}

Before turning to the onset analysis, I examine the descriptive properties
of $E^G$ and $E^O$ to establish that they behave as the theory predicts.

The dyad-level predictions reveal a structural asymmetry between the two
sides of intervention.  The highest predicted probabilities of
government-biased intervention involve superpowers projecting influence into
client states---the United States supporting the Philippines in 1972
($\hat{p}_\text{gov} = 0.81$), Laos in 1968, Cambodia in 1972.  For
opposition-biased intervention, the picture shifts: the maximum is Somalia
opposing the Ethiopian government in 1980 ($\hat{p}_\text{opp} = 0.77$), followed
closely by Bulgaria targeting Greece in 1947 and Libya targeting Chad in 1984,
and all ten of the highest-ranked opposition predictions involve non-P5
interveners.  A measurement framework restricted to P5 states
would miss the majority of the opposition-biased shadow.

Table~\ref{tab:shadow-dyads} makes this concrete for three canonical onsets, reporting
the five states to which the classifier assigns the highest government- and
opposition-biased intervention probabilities alongside an indicator of whether
the state actually intervened in that direction.

\begin{table}[t]
  \centering
  \small
  \setlength{\tabcolsep}{4pt}
  \begin{tabular}{@{}lr c @{\quad} lr c@{}}
    \toprule
    \textbf{Gov-biased} & $\hat{p}$ & & \textbf{Opp-biased} & $\hat{p}$ & \\
    \midrule
    \multicolumn{6}{c}{\textbf{Angola (1976) --- 5 actual interventions}} \\
    \cmidrule(lr){1-6}
    USSR$^\ast$          & .20 & $\checkmark$ & Zaire              & .49 &              \\
    Cuba                 & .18 & $\checkmark$ & South Africa       & .19 & $\checkmark$ \\
    United States$^\ast$ & .15 &              & United States$^\ast$ & .04 & $\checkmark$ \\
    Zambia               & .07 &              & Portugal           & .03 &              \\
    Portugal             & .05 &              & USSR$^\ast$        & .02 &              \\
    \addlinespace
    \multicolumn{6}{c}{\textbf{Ethiopia (1975) --- 3 actual interventions}} \\
    \cmidrule(lr){1-6}
    USSR$^\ast$          & .14 & $\checkmark$ & Somalia            & .66 & $\checkmark$ \\
    South Yemen          & .12 &              & United States$^\ast$ & .06 &              \\
    United States$^\ast$ & .11 &              & North Yemen        & .06 &              \\
    Iran                 & .04 &              & Sudan              & .04 &              \\
    Saudi Arabia         & .03 &              & South Yemen        & .02 &              \\
    \addlinespace
    \multicolumn{6}{c}{\textbf{Afghanistan (1978) --- 4 actual interventions}} \\
    \cmidrule(lr){1-6}
    USSR$^\ast$            & .29 & $\checkmark$ & Pakistan             & .45 & $\checkmark$ \\
    United Kingdom$^\ast$  & .11 &              & United States$^\ast$ & .08 & $\checkmark$ \\
    China$^\ast$           & .09 &              & United Kingdom$^\ast$  & .06 &              \\
    United States$^\ast$   & .07 &              & China$^\ast$         & .05 &              \\
    Iran                   & .07 &              & Iran                 & .04 & $\checkmark$ \\
    \bottomrule
  \end{tabular}
  \caption{Top five predicted interveners for three canonical onsets.
    Out-of-fold probabilities, averaged across 25 imputation draws.
    $^\ast$P5 member. $\checkmark$ indicates the state actually
    intervened in that direction (Regan coding).  Ten of 12 actual
    interventions appear in the top five on the correct side.}
  \label{tab:shadow-dyads}
\end{table}

Across the three cases, 10 of 12 actual interventions appear in the top five
on the correct side---and in each case, restricting the measure to P5 states
would discard the single most important opposition-biased intervener: Zaire
in Angola, Somalia in Ethiopia, Pakistan in Afghanistan.\footnote{%
  The United States appears in the top five government-biased predictions for
  all three cases, yet actually intervened on the opposition side in Angola and
  Afghanistan.  The model assigns substantial probability mass to \textit{both} sides
  but defaults toward government-biased intervention because that is the modal
  US posture across the full sample.  This directional ambiguity for superpowers
  is one reason the country-year aggregation is preferable to raw dyad-level
  predictions: summing across all potential interveners preserves the correct
  signal that these countries faced unusually high total intervention pressure,
  even where the classifier hedges on which side a specific superpower would take.}

Figure~\ref{fig:shadow-kde} shows the cross-sectional distribution of $E^G$ and $E^O$
across all country-years.  Both measures are heavily right-skewed: the
median country-year faces modest intervention expectations, but a long
right tail captures the handful of states embedded in dense alliance
networks and superpower competition.

\begin{figure}[t]
  \centering
  \includegraphics[width=0.85\textwidth]{fig-shadow-kde.pdf}
  \caption{Kernel density estimates of $E^G$ and $E^O$ across all
    country-years (1946--2014), averaged across 25 imputation draws.}
  \label{fig:shadow-kde}
\end{figure}

Figure~\ref{fig:shadow-ts} illustrates the temporal dynamics for three countries
with well-documented intervention environments.  South Vietnam shows the
cleanest signal: both measures ramp from approximately 0.6 in 1960 to
$E^G = 6.68$ and $E^O = 5.21$ at the peak of the war in 1972---an order
of magnitude above typical countries.
Afghanistan illustrates a regime-switching dynamic: the 1978 Saur Revolution
raises $E^G$ to 0.84, but by 1979 $E^O$ overtakes $E^G$ ($1.17$ vs $0.85$),
correctly capturing the surge in mujahideen support from Pakistan, the
United States, and others.
Egypt provides a critical out-of-sample check: it experiences no civil war
onset in our data, so the shadow measures are pure counterfactual predictions.
The classifier nonetheless tracks the geopolitical environment---the 1967
Six Day War spike ($E^G = 7.67$), the post-1972 Soviet expulsion drop, the
1991 Gulf War spike---evidence of genuine temporal tracking, not mere
onset-fitting.

\begin{figure}[t]
  \centering
  \includegraphics[width=0.85\textwidth]{fig-shadow-ts.pdf}
  \caption{Shadow measure time series ($E^G$ and $E^O$) for five countries
    with well-documented intervention environments.  Probabilities averaged
    across 25 imputation draws.}
  \label{fig:shadow-ts}
\end{figure}

$E^G$ correlates positively with both existing measures of anticipated
intervention: the Lake security hierarchy used by \citet{cunningham2016}
($r = 0.29$, $n = 161$) and the structural P5 intervention probabilities
estimated by \citet{gibilisco2022} ($r = 0.53$, $n = 145$).  The stronger
correlation with G\&M is expected---both capture intervention propensity
directly---while the moderate hierarchy correlation reflects the fact that
US dominance is only one component of the broader intervention environment
(Appendix~\ref{app:cunningham}, Figure~\ref{fig:hierarchy}).
